% Источники: exam/materials/преподаватель/2020/20-04-2020-Формализация_математической_модели___пример_к_лекции_9_.pdf; exam/materials/моделирование.pdf, разделы 14 и 48.
\section{Методика построения аналитической динамической модели.}

\textbf{Динамическая модель} отражает изменение состояния исследуемой системы или объекта в заданный или формируемый период времени. Пример динамической модели — баллистическая модель.

\subsection*{Методика построения (7 шагов)}

\begin{enumerate}
  \item \textbf{Оценивается постановка задачи} и процесс изменения объекта во времени: строится графическая схема в начальный момент $t_0$, в некоторый фиксированный $t^*$ и в конечный момент $T$. Для описания процесса может быть также построена диаграмма в выбранной графической нотации.

  \item По схеме \textbf{обозначаются начальные и конечные входные данные}; выделяются переменные, которые являются выходными на одном этапе и входными на другом; определяются выходные данные, полученные в процессе моделирования.

  \item \textbf{Определяются физико-химические воздействия} на объект, влияющие на изменения выходных данных.

  \item \textbf{Оценивается возможность упрощения модели} — ряд воздействий не учитывается.

  \item \textbf{Проводится формализация модели:} с учётом известных формул описываются учитываемые воздействия и составляется общее уравнение (равенство, неравенство, система уравнений и т.п.), связывающее входные и выходные данные (в левой и правой части уравнения — величины одинаковой размерности).

  \item \textbf{Конкретизируется описание} каждого отдельного воздействия.

  \item \textbf{Повторно оценивается возможность упрощения модели:} число неизвестных не должно превышать числа уравнений в системе; уточняются значения коэффициентов и констант (в т.ч. начальных и конечных значений переменных), входящих в описание.
\end{enumerate}

\clearpage
